We begin by defining the pointed topologically enriched category $\J$, which serves as the indexing category for chain complexes in the topologically enriched setting.

\begin{defn}
We define the pointed topologically enriched category $\J$ as follows:
\begin{itemize}
    \item The objects of $\J$ are the integers together with a distinguished object $\infty$:
    \begin{equation*}
        \Ob(\J) = \mathbb{Z} \cup \{ \infty \}.
    \end{equation*}
    \item The morphism spaces $\Hom_{\J}(j, i)$ are defined for $i, j \in \mathbb{Z} \cup \{ \infty \}$ by:
    \begin{equation*}
        \Hom_{\J}(j, i) = 
        \begin{cases}
            \{ \id, \infty \} & \text{if } i = j, \\
            \{ d_i, \infty \} & \text{if } j = i + 1, \\
            (\Hlf^{J_{j,i}})_+ & \text{if } j > i + 1, \\
            \{ \infty \} & \text{if } i = \infty \text{ or } j = \infty,
        \end{cases}
    \end{equation*}
    where $X_+$ denotes the one-point compactification of a space $X$, and $\Hlf$ is the half-line $[0, \infty)$. The set $J_{j,i}$ is defined by:
    \begin{equation}\label{eq:J_indices}
        J_{j,i} = \{ \ell \in \mathbb{Z} \mid j > \ell > i \}.
    \end{equation}
    The spaces $\Hom_{\J}(j, i)$ are pointed at the point $\infty$.
    \item The maps
    \begin{equation*}
        J_{k,j}\sqcup J_{j,i} \hookrightarrow J_{k,i}
    \end{equation*}
    induce a map
    \begin{equation*}
        \Hlf^{J_{k,j}}\times \Hlf^{J_{j,i}} \hookrightarrow \Hlf^{J_{k,i}}.
    \end{equation*}
    By taking $(-)_+$ we obtain the composition in $\J$:
    \begin{equation*}
        (\Hlf^{J_{k,j}})_+\smashSpaces (\Hlf^{J_{j,i}})_+ \to (\Hlf^{J_{k,i}})_+.
    \end{equation*}
\end{itemize}
\end{defn}

\begin{rem}
Intuitively, the morphisms in $\J$ correspond to the differentials and higher compositions in a chain complex. The space $\Hlf^{J_{j,i}}$ models the higher nullhomotopies between compositions of differentials.
\end{rem}\

\begin{defn}
    Let $C$ be a pointed topologically enriched category. 
    We define (topological, coherent) chain complexes in $C$ as strict functors $\Fun_+(\J,C)$.
\end{defn}

Let $Sp^O$ be the category of orthogonal spectra. This is a topologically enriched model category for the stable homotopy category, as developed in \cite{MMSS01}. It is equipped with a symmetric monoidal smash product $\otimes$ that gives a Quillen equivalent model to the $\infty$-category of spectra.
In this model, the suspension spectrum functor 
\begin{equation*}
    \Sigma^{\infty}: (Top_{+}, \wedge, S^{0}) \to (Sp^{O}, \otimes, \mathbb{S})
\end{equation*}
is strongly symmetric monoidal and preserves coproducts.

\begin{constr}
    Let $\fC$ be a flow category with neat embedding $e$ and normal framing $\varphi$. We define $K_\fC \in \Fun_+(\J,\Sp^O)$ by
    \[
        K_\fC(i) = \WedgeSp_{\norm{x}=i} \bbS
    \]
    To define the morphisms it's suffice to we will construct maps
    \[
        \Hom_\J(j,i) \to \Hom_{\Top_+}(S^{N(j-i)}, S^{N(j-i)})
    \]
    and then by taking $\Sinf$ and desuspending we get a map
    \[
        \Hom_\J(j,i) \to \Hom_{\Sp^O}(\bbS, \bbS)
    \]
     we use the $(\smashSpaces, \Hom)$ adjunction. It suffices to provide maps:
    \[
        g_{y,x}:\Hom_\J(j,i)\smashSp \bbS \to \bbS
    \]
    for each pair $y,x \in \fC$ with indices $j,i$.
    Let $y$ and $x$ be as above; we define the maps
    \[
        S^{N(\norm{y}-\norm{x})}\smashSp (\Hlf^{J_{y,x}})_+ \to  Th(e_{y,x}) \to S^{N(\norm{y}-\norm{x})},
    \]
    where the first map is the collapse map of $e_{y,x}$ and the second is (the compactification of) the projection of $\varphi_{y,x}$ (see \cref{eq:ThomMap}). By taking $\Sinf$ and desuspending this map, we obtain:
    \[
        g_{y,x}: \bbS \smashSp (\Hlf^{J_{y,x}})_+ \to  Th(e_{y,x}) \to \bbS.    
    \]
    Functoriality of $K_\fC$ amounts to the equation 
    \[
        \partial_y g_{z,x} = g_{y,x} \circ g_{z,y},
    \]
    which holds due to the compatibility of the neat embedding and the framing (see \cref{eq:ThomFunctoriality}).
\end{constr}


We call a functor $\Fun_+(\J,\Top_+)$ a \textbf{topological chain complex}. Our next goal is to start with an embedded framed flow category and produce such a topological chain complex. The idea is to use the Pontryagin-Thom collapse map (\cref{eq:ThomMap}) in a compatible way to "collapse" the flow category and then use $\J$ to encode the maps between spheres and their higher homotopies.

\begin{constr}\label{con:topChOfCat}
    Let $\fC$ be a flow category with a neat embedding $e$ and normal framing $\varphi$.
    Assume\todo{why we need this? how CK of CJS do this construction?} that $\fC$ has bounded index 
\begin{equation*}
    b \geq \norm{-} \geq a.
\end{equation*}
Without loss of generality, we take $a=0$. We define a pointed topologically enriched functor $K_{\fC} \in \Fun_+(\J,\Top_+)$ as follows.
    
    First, we set
    \begin{equation*}
        K_{\fC}(i) = \WedgeSpaces_{\norm{x}=i} S^{Nb}.
    \end{equation*}
    To define the morphisms, we use the $(\smashSpaces, \Hom)$ adjunction. It suffices to provide maps:
    \begin{equation*}
        g_{y,x}:\Hom_{\J}(j,i)\smashSpaces S^{Nb} \to S^{Nb}
    \end{equation*}
    for each pair $y,x \in \fC$ with indices $j,i$.
    Let $y$ and $x$ be as above; we define the maps
    \begin{equation*}
        S^{N(\norm{y}-\norm{x})}\smashSpaces (\Hlf^{J_{y,x}})_+ \to  Th(e_{y,x}) \to S^{N(\norm{y}-\norm{x})},
    \end{equation*}
    where the first map is the collapse map of $e_{y,x}$ and the second is (the compactification of) the projection of $\varphi_{y,x}$ (see \cref{eq:ThomMap}). By suspending this map, we obtain:
    \begin{equation*}
        g_{y,x}: S^{Nb}\smashSpaces (\Hlf^{J_{y,x}})_+ \to  Th(e_{y,x}) \to S^{Nb}.    
    \end{equation*}
    Functoriality of $K_{\fC}$ amounts to the equation 
    \begin{equation*}
        \partial_y g_{z,x} = g_{y,x} \circ g_{z,y},
    \end{equation*}
    which holds due to the compatibility of the neat embedding and the framing (see \cref{eq:ThomFunctoriality}).
\end{constr}

For a flow category $\fC$, we define $\fC^{[b,a]}$ to be the full subcategory generated by objects $x \in \fC$ satisfying the index condition
\begin{equation*}
    b \geq \norm{x} \geq a.
\end{equation*}
We denote by $K_{[b,a]}$ the topological chain complex of $\fC^{[b,a]}$. Note that there is a natural transformation 
\begin{equation*}
    \beta_{b+1}:K_{[b,a]} \to K_{[b+1,a]}
\end{equation*}
induced by suspending further to $N(b+1)$; if $i=b+1$, we take $\beta_{b+1}(i)$ to be the zero map. Similarly, we have maps 
\begin{equation*}
    \alpha_{a-1}:K_{[b,a-1]} \to K_{[b,a]}.
\end{equation*}

\todo{Do the following def works? do we get a functor from $\J$?, see also the last todo}
\begin{defn}\label{def:unbounded_K}
    We extend \cref{con:topChOfCat} to unbounded flow categories by taking the limit and colimit:
    \begin{equation*}
        K_{\fC} = \lim_a \colim_b K_{[b,a]}.
    \end{equation*}
\end{defn}

\begin{constr}\label{con:SpChOfCat}
    Let $\fC$ be an embedded framed flow category bounded by $b$ and $0$. Using \cref{con:topChOfCat} and post-composition with $\Sinf$, we construct a functor $K_{\fC}' \in \Fun(J,\SPO)$ with 
    \begin{equation*}
        K_{\fC}'(i) = \WedgeSp_{\norm{x}=i} \bbS^{Nb}.
    \end{equation*}
    Using post-composition with $\Sigma^{-Nb}: \Sp \to \Sp$, we define the chain complex of $\fC$ to be 
    \begin{equation*}
        K_{\fC} \coloneqq \Sigma^{-Nb} K_{\fC}'
    \end{equation*}
    which has $\bbS$ in each level.
\end{constr}

\begin{prop}\label{prop:J_equivalence}
There is an equivalence of pointed $\infty$-categories $N(\Ch) \simeq N_{\Top}(\J)$, where $\Ch$ is the indexing category of chain complexes (see \cref{def:Ch}) and $N, N_{\Top}$ are the nerve and the topological nerve, respectively.
\end{prop}

\begin{proof}
The morphism spaces $\Hom_{\J}(j,i)$ are contractible when $j > i + 1$ (since $(\Hlf^{J_{j,i}})_+$ is contractible). Furthermore, we have the following identifications:
\begin{equation*}
    \{ \id, \infty \} \simeq S^0 \text{ for } i = j, \quad \text{and} \quad \{ d_i, \infty \} \simeq S^0 \text{ for } j = i + 1.
\end{equation*}
The equivalence follows by considering both $\Ch$ and $\J$ as models for the same $\infty$-category.
\end{proof}

\begin{prop}\label{prop:FF_are_Ch}
    In light of this equivalence, we can start with a framed flow category and obtain a coherent chain complex in spectra:
    \begin{equation*}
    \Fun_{\Top}(\J,\SPO) \simeq \Fun(N_{\Top}(\J),N_{\Top}(\SPO)) \simeq \Fun(N(\Ch),\Sp) \eqqcolon \Ch(Sp).
\end{equation*}
\end{prop}

By identifying $1$-categories with their nerves, we end this section by motivating the next. Let $\fC$ be a bounded, embedded, framed flow category and let $K \in \Fun(\Ch,\Sp)$ be its coherent chain complex. By \cref{sec:Ch}, we obtain a spectral sequence $(\rnp{E}, d_r)$ where 
\begin{equation*}
    \pg{E}{1}{n}{p} = \pi_{n-p} K_p = \bigoplus_{\norm{x}=p} \pi_{n-p}\bbS.
\end{equation*}
Identifying stable homotopy groups with framed bordism via the isomorphism
\begin{equation*}
    \pi_{n-p} \bbS \simeq \Omega^{\fr}_{n-p},
\end{equation*}
we observe that for $y,x \in \fC$ with $\norm{y}=p$ and $\norm{x}=p-1$, the first differential 
\begin{equation*}
    d_1\colon \Omega^{\fr}_{n-p}\<{y} \to \Omega^{\fr}_{n-p}\<{x}
\end{equation*}
is given by the formula
\begin{equation*}
    d_1(N) = N \times M_{y,x},
\end{equation*}
where $M_{y,x}$ is the morphism manifold in $\fC$ and $N$ is a representative for its bordism class. Survival to the second page corresponds to the existence of a flow category with objects $s,y,x$ and morphism manifolds 
\begin{equation*}
    M_{y,x}, \quad N=M_{s,y}, \quad \text{and} \quad W=M_{s,x}
\end{equation*}
such that their boundaries satisfy the relation
\begin{equation*}
    \partial W = M_{y,x} \times N.
\end{equation*}
